\pdfvariable suppressoptionalinfo 611
\providecommand{\CRevisionRound}{2}
\documentclass[9pt]{extarticle}
\usepackage[a4paper,margin=0.64in]{geometry}
\usepackage{fontspec}
\setmainfont{Latin Modern Roman}
\setsansfont{Latin Modern Sans}
\setmonofont{Latin Modern Mono}
\usepackage[english]{babel}
\babelprovide[import]{chinese-simplified}
\babelfont[chinese-simplified]{rm}[Renderer=Harfbuzz,Language=Default]{Droid Sans Fallback}
\usepackage{amsmath,amssymb,booktabs}
\emergencystretch=1em
\title{Relaxed Douglas--Rachford Dynamics for Two Subspaces:\\
An All-Parameter Principal-Angle Atlas}
\author{Route-A structural certificate C197}
\date{}
\begin{document}
\maketitle
\begin{abstract}
For every pair of linear subspaces in finite-dimensional real Hilbert space
and every real relaxation parameter, we give the complete orthogonal block
atlas of the relaxed Douglas--Rachford map.  It yields the fixed space, sharp
convergence window and rate, optimal relaxation, shadow limit, endpoint
rotation dynamics, instability boundary, all power traces, and a closed finite
determinant.  Exact rational sentinels are reconstructed independently from
projections and by symbolic algebra; they test conventions without replacing
the all-subspace proof.  The source-native orthogonal endpoint is a formal
operator hint, but projection geometry has no intrinsic prime carrier or
logarithmic clock, so the candidate fails Route A.
\end{abstract}
\noindent\textbf{Keywords:} Douglas--Rachford splitting; principal angles;
relaxed projections; convergence rate; finite determinant; Route-A obstruction.

\begin{otherlanguage}{chinese-simplified}
\begin{center}{\large 中文摘要}\end{center}
本文对有限维实希尔伯特空间中任意两个线性子空间及任意实松弛参数，给出松弛
\foreignlanguage{english}{Douglas--Rachford}映射的完整主角分块。该分块一次确定不动空间、
收敛区间与精确速率、最优松弛、影子极限、端点旋转、失稳边界、全部幂迹与有限行列式。
有理角度证书仅检验符号与约定，不替代全参数证明。正交端点虽然是同钟自然算子，却没有
内生素数载体或对数时钟，故不能通过\foreignlanguage{english}{Route A}。

\noindent 关键词：投影算法；主角；松弛动力学；收敛速率；算术阻碍。
\end{otherlanguage}

\section{Frozen map and canonical planes}
Let $U,V\subset H$ be linear subspaces, $P_W$ the orthogonal projection and
$R_W=2P_W-I$.  One algorithmic tick is
\[
 T_\lambda=(1-\lambda)I+\frac{\lambda}{2}(I+R_VR_U),\qquad \lambda\in\mathbb R.
\]
This is the two-subspace specialization of the classical splitting lineage
of Douglas and Rachford~\cite{DR56}; the sharp subspace convergence location
and Friedrichs-angle rate are due to Bauschke et al.~\cite{B14}.

The cosine--sine decomposition splits $H$ orthogonally into
\[
 F=(U\cap V)\oplus(U^\perp\cap V^\perp),\qquad
 M=(U\cap V^\perp)\oplus(U^\perp\cap V),
\]
and principal planes $E_j$ with $0<\theta_j<\pi/2$.  On $F$ and $M$ the map is
$I$ and $(1-\lambda)I$, respectively.  Choosing
$U\cap E_j=\operatorname{span}(e_1)$ and
$V\cap E_j=\operatorname{span}(\cos\theta_j e_1+\sin\theta_j e_2)$ gives
\begin{equation}\label{eq:block}
 T_\lambda|_{E_j}=\begin{pmatrix}
 1-\lambda\sin^2\theta_j&-\lambda\sin\theta_j\cos\theta_j\\
 \lambda\sin\theta_j\cos\theta_j&1-\lambda\sin^2\theta_j
 \end{pmatrix}.
\end{equation}
Its conjugate eigenvalues have squared modulus
\begin{equation}\label{eq:rho}
 \rho_j^2=1-\lambda(2-\lambda)\sin^2\theta_j.
\end{equation}

\section{Complete dynamics}
For $\lambda\ne0$, $\operatorname{Fix}T_\lambda=F$; at $\lambda=0$ the map is
the identity.  Equations~\eqref{eq:block}--\eqref{eq:rho} and orthogonality of
the decomposition prove the following exhaustive classification.
\begin{center}
\begin{tabular}{@{}lll@{}}\toprule
parameter & off-fixed behavior & conclusion\\ \midrule
$\lambda=0$ & identity & every point fixed\\
$0<\lambda<2$ & strict contraction & $T_\lambda^n\to P_F$\\
$\lambda=2$ & signs and plane rotations & orthogonal recurrence boundary\\
$\lambda\notin[0,2]$ & modulus $>1$ if nontrivial & unstable off $F$\\ \bottomrule
\end{tabular}
\end{center}
In the contracting window the exact operator-norm rate is the largest existing
quantity among
\[
 |1-\lambda|,\qquad
 \sqrt{1-\lambda(2-\lambda)\sin^2\theta_j}.
\]
Thus $\lambda=1$ uniquely minimizes every nontrivial block simultaneously; its
generic rate is $\cos\theta_F$.  Moreover
$P_UT_\lambda^n x\to P_{U\cap V}x$, because $P_UP_F=P_{U\cap V}$.

At $\lambda=2$, $T_2=R_VR_U$: it is $I$ on $F$, $-I$ on $M$, and rotation by
$2\theta_j$ on $E_j$.  It has finite order exactly when every
$\theta_j/\pi$ is rational (with the mismatch signs included in the least
common order).  Hence the endpoint is not silently counted as convergence.

\ifnum\CRevisionRound>0
\section{Trace and determinant closure}
Put $f=\dim F$, $h=\dim M$,
$a_j=1-\lambda\sin^2\theta_j$, and $d_j=\rho_j^2$.  Then the complete finite
factorization is
\begin{equation}\label{eq:det}
 \det(I-zT_\lambda)=(1-z)^f(1-(1-\lambda)z)^h
 \prod_j(1-2a_jz+d_jz^2).
\end{equation}
The contribution $\tau_{j,n}$ of a principal plane to
$\operatorname{tr}(T_\lambda^n)$ obeys
$\tau_{j,0}=2$, $\tau_{j,1}=2a_j$ and
$\tau_{j,n}=2a_j\tau_{j,n-1}-d_j\tau_{j,n-2}$.  Thus the same block theorem
closes every trace, not only the asymptotic rate.  Formula~\eqref{eq:det} is a
finite algorithmic determinant; no periodic-orbit or target ownership is
introduced by notation.
\fi

\ifnum\CRevisionRound>1
\section{Exact certificate and strict Route-A stop}
The release ledger contains 28 rational generic blocks and 21 composite
eight-dimensional models: 112 matrix cells and 441 power-trace cells plus all
determinant coefficients.  The producer evaluates~\eqref{eq:block}; an
independent checker instead constructs $P_U,P_V,R_U,R_V$ and takes direct
matrix powers.  A separate symbolic derivation reconstructs all rows.  Replay
is byte exact, and twelve repaired-hash semantic attacks plus one stale-hash
attack are rejected.  The rational sentinels are regression, not proof of the
cosine--sine decomposition.

Projection geometry has no intrinsic rational-prime primitive carrier,
prime-power repetition or $\log p$ clock.  Prime and composite dimensions obey
the same theorem.  The finite determinant has neither target divisor nor
functional equation.  Although $T_2$ is a source-native same-clock orthogonal
map, relabeling it as a quantum operator would be post hoc.  Therefore
\[
 (A0,A1,A2,A3,A4)=(\mathrm{FAIL},\mathrm{FAIL},\mathrm{FAIL},
 \mathrm{FAIL},\mathrm{FORMAL\ HINT}),
\]
overall \texttt{ROUTE\_A\_REJECTED}, Route B false, under
\texttt{NO\_BAD\_EULER\_OR\_ROOT\_NUMBER}.

\paragraph{Scope firewall.}
We claim neither priority for splitting or its angle rate, nor an
infinite-dimensional norm theorem without closed-sum hypotheses.  No target
zero or prime table, arithmetic local datum, Euler factor, root number,
automorphy, target analytic structure, Hilbert--P\'olya operator, external peer
review or acceptance score is present.
\fi

\ifcase\CRevisionRound
\paragraph{Revision focus.} Round 0 freezes the operator and proves the
principal-plane atlas and complete relaxation classification.
\or
\paragraph{Revision focus.} Round 1 adds fixed/mismatch boundaries, sharp
optimality, the shadow theorem and full trace--determinant closure.
\or
\paragraph{Revision focus.} Round 2 adds source ownership, independent exact
validation, endpoint periodicity and the strict Route-A claim boundary.
\fi

\section*{Declarations}
\paragraph{Data and code availability.} Package-local exact evidence and
deterministic code accompany this manuscript.
\paragraph{Ethics.} No human, animal, clinical, personal or sensitive data are
used; approval is not applicable.
\paragraph{Author contributions (CRediT).} This anonymous certificate records
Conceptualization, Formal analysis, Software, Validation, Writing---original
draft, and Writing---review and editing. AI systems are not authors.
\paragraph{Funding.} No external funding is reported.
\paragraph{Conflicts of interest.} None known.
\paragraph{AI-use disclosure.} An AI coding assistant supported drafting and
exact-code development; it was not an external reviewer or independent error
process.

\begin{thebibliography}{2}
\bibitem{DR56} J. Douglas Jr. and H. H. Rachford Jr., ``On the numerical
solution of heat conduction problems in two and three space variables,''
\emph{Trans. Amer. Math. Soc.} 82 (1956), 421--439.
DOI: 10.1090/S0002-9947-1956-0084194-4.
\bibitem{B14} H. H. Bauschke, J. Y. Bello Cruz, T. T. A. Nghia, H. M. Phan
and X. Wang, ``The rate of linear convergence of the Douglas--Rachford
algorithm for subspaces is the cosine of the Friedrichs angle,''
\emph{J. Approx. Theory} 185 (2014), 63--79.
DOI: 10.1016/j.jat.2014.06.002.
\end{thebibliography}
\end{document}
